\section{Table spans}

Carve tables carry cell spans with two markers: \texttt{\textless{}}
merges a cell into its left neighbor (colspan) and \texttt{\^{}} merges
a cell into the one above it (rowspan). pandoc-carve rewrites these into
Pandoc's origin-cell span model (the origin cell holds the width/height;
covered cells are omitted), so the merged regions survive into LaTeX
\texttt{\textbackslash{}multicolumn}/\texttt{\textbackslash{}multirow},
DOCX merged cells, Typst, and the rest - not just HTML.

\begin{longtable}[]{@{}lll@{}}
\caption{Sales by region (North 2024 spans both years; South spans two
rows)}\tabularnewline
\toprule\noalign{}
Region & 2024 & 2025 \\
\midrule\noalign{}
\endfirsthead
\toprule\noalign{}
Region & 2024 & 2025 \\
\midrule\noalign{}
\endhead
\bottomrule\noalign{}
\endlastfoot
North & \multicolumn{2}{l@{}}{%
100} \\
\multirow{2}{=}{South} & 40 & 60 \\
& 30 & 50 \\
\end{longtable}
